\chapter{Model, notation and assumptions}
\label{chap:model}

This chapter sets up the probabilistic model of a sequential multi-treatment
experiment, the estimators used, and the family of allocation rules
($\aRTS$) whose asymptotic behaviour is analysed in the rest of the blueprint.
It follows Section~2 and Section~3 of the paper.

\textbf{Lean remark.}
Throughout, everything lives on a fixed probability space $(\Omega, \filt, \Prob)$.
We work with $K \geq 2$ treatments (``arms''), indexed by $[K] := \{1, \dots, K\}$.
Random objects are indexed by a time $m \in \N$ (the patient number).
A convenient Lean encoding is to fix once and for all:
\begin{itemize}
  \item a natural number $K$ with $2 \le K$,
  \item a probability space $(\Omega, \filt, \Prob)$,
  \item the assignment process $X : \N \to \Omega \to (\text{Fin } K \to \{0,1\})$,
  \item the response process $\xi : \N \to \text{Fin } K \to \Omega \to \R$.
\end{itemize}

\section{The sequential experiment}
\label{sec:model_experiment}

\begin{definition}[Assignment vector]\label{def:assignment}
\lean{Learning.actionIndicator}\leanok
For each patient $m \ge 1$, the \emph{assignment vector}
$X_m = (X_{m,1}, \dots, X_{m,K})$ takes values in $\{0,1\}^K$ and satisfies
$\sum_{k=1}^K X_{m,k} = 1$. The event $\{X_{m,k} = 1\}$ means that patient $m$ is
assigned to arm $k$.
\end{definition}

\begin{definition}[Responses and parameters]\label{def:responses}
For each arm $k \in [K]$, $\{\xi_{m,k} : m \ge 1\}$ is a sequence of i.i.d.\
real random variables, the potential responses of patients on treatment $k$.
The parameter of interest for arm $k$ is the mean
$\theta_k := \E[\xi_{1,k}]$, and $\Params = (\theta_1, \dots, \theta_K) \in \R^K$.
We write $V_k := \Var[\xi_{1,k}]$ for the variance of the response of arm $k$.
\end{definition}

\begin{definition}[Allocation counts]\label{def:counts}
\lean{AlphaRAR.count}\leanok
\uses{def:assignment}
The number of patients assigned to arm $k$ up to time $n$ is
\[
  N_{n,k} := \sum_{j=1}^n X_{j,k}, \qquad N_n := (N_{n,1}, \dots, N_{n,K}).
\]
The vector $N_n / n$ is the \emph{observed allocation proportion} at time $n$.
\end{definition}

\begin{lemma}[Counts sum to time]\label{lem:counts_sum}
\lean{AlphaRAR.counts_sum}\leanok
\uses{def:counts}
For all $n \ge 1$, $\sum_{k=1}^K N_{n,k} = n$, and hence
$\sum_{k=1}^K \frac{N_{n,k}}{n} = 1$.
\end{lemma}

\begin{proof}\leanok
\uses{def:assignment}
Immediate from $\sum_{k=1}^K X_{j,k} = 1$ and summing over $j \le n$.
\end{proof}

\begin{definition}[Filtration]\label{def:filtration}
\lean{Learning.IsAlgEnvSeq.filtration}\leanok
\uses{def:assignment, def:responses}
$\filt_m := \sigma\big((X_i, \xi_i)_{1 \le i \le m}\big)$ is the filtration
generated by the assignments and observed outcomes of the first $m$ patients,
with $\filt_0$ the trivial $\sigma$-algebra.
\end{definition}

\begin{definition}[RAR procedure]\label{def:rar}
\uses{def:filtration}
A \emph{response-adaptive randomization (RAR) procedure} is a sequence
$(p_m)_{m \ge 1}$ of $\filt_m$-measurable random vectors taking values in the
probability simplex
\[
  \simplex_K := \Big\{ p \in [0,1]^K : \sum_{k=1}^K p_k = 1 \Big\},
\]
such that the assignment $X_{m+1}$ is drawn conditionally on $\filt_m$ according
to $p_m$:
\[
  \Prob(X_{m+1} = k \mid \filt_m) = p_{m,k}, \qquad k \in [K].
\]
\end{definition}

\textbf{Lean remark.} The conditional-law requirement is best captured through
the conditional expectation $\E[X_{m+1,k} \mid \filt_m] = p_{m,k}$
(see Definition~\ref{def:cond_prob}), which is the only consequence used in the proofs.

\begin{definition}[Target function and target allocation]\label{def:target}
A \emph{target function} is a map $\Target : \hilb \to \simplex_K$ defined on an
open domain $\hilb \subseteq \R^K$, written
$\Target(z) = (\Target_1(z), \dots, \Target_K(z))$ with
$\sum_{k=1}^K \Target_k(z) = 1$. The \emph{target allocation} is
$v := \Target(\Params)$. A RAR procedure is an \emph{adaptive targeting
mechanism} if $N_n/n \to v$.
\end{definition}

\section{Sequential estimation}
\label{sec:model_estimation}

\begin{definition}[Sequential estimator]\label{def:estimator}
\lean{AlphaRAR.estimator}\leanok
\uses{def:counts, def:responses}
Given initial values $\theta_{0,k} \in \R$, the sequential estimator of
$\theta_k$ at stage $m$ is
\[
  \hat{\theta}_{m,k} := \frac{\sum_{j=1}^m X_{j,k}\, \xi_{j,k} + \theta_{0,k}}{N_{m,k} + 1}.
\]
We write $\estParams_m := (\hat{\theta}_{m,1}, \dots, \hat{\theta}_{m,K})$.
\end{definition}

\begin{definition}[Plug-in target]\label{def:plugin_target}
\lean{AlphaRAR.aRTSTarget}\leanok
\uses{def:estimator, def:target}
The plug-in estimator of the target allocation is
$\estTarget_m := \Target(\estParams_m)$, with components $\hat{\rho}_{m,k}$.
It is well defined once $\estParams_m \in \hilb$, which is guaranteed by
Condition~\ref{def:condB}.
\end{definition}

\begin{definition}[Bahadur representation]\label{def:bahadur}
\uses{def:estimator}
An estimator $\hat{\theta}_{m,k}$ of $\theta_k$ admits a \emph{Bahadur
representation} if
\[
  \hat{\theta}_{m,k} = \frac{1}{N_{m,k}}\sum_{j=1}^m X_{j,k}\, \xi_{j,k}
    + o\!\big(N_{m,k}^{-1/2}\big), \qquad m \to \infty.
\]
\end{definition}

\begin{lemma}[The sequential estimator is Bahadur]\label{lem:estimator_bahadur}
\lean{AlphaRAR.estimator_sub_eq}\leanok
\uses{def:estimator, def:bahadur, def:Q}
On the event $\{N_{m,k} \to \infty\}$, the estimator of
Definition~\ref{def:estimator} admits the Bahadur representation of
Definition~\ref{def:bahadur}. The formalized statement \texttt{estimator\_sub\_eq} is the sharper
\emph{exact} error identity
\[
  \hat{\theta}_{m,k} - \theta_k = \frac{Q_{m,k} + (\theta_{0,k} - \theta_k)}{N_{m,k} + 1},
\]
which holds pathwise for every $m$ (with $N_{m,k}+1 \ne 0$, automatic for $\{0,1\}$-indicators) and
from which the $o(N^{-1/2})$ Bahadur form is immediate.
\end{lemma}

\begin{proof}\leanok
The numerator of $\hat{\theta}_{m,k} - \theta_k = \frac{\sum_j X_{j,k}\xi_{j,k} + \theta_{0,k}}{N_{m,k}
+ 1} - \theta_k$ is $\sum_j X_{j,k}\xi_{j,k} + \theta_{0,k} - \theta_k(N_{m,k}+1) = (\sum_j
X_{j,k}\xi_{j,k} - \theta_k N_{m,k}) + (\theta_{0,k} - \theta_k) = Q_{m,k} + (\theta_{0,k} -
\theta_k)$, giving the exact identity. On $\{N_{m,k}\to\infty\}$ the offset
$\frac{\theta_{0,k}-\theta_k}{N_{m,k}+1} = O(N_{m,k}^{-1}) = o(N_{m,k}^{-1/2})$, which is the Bahadur
form.
\end{proof}

\textbf{Lean remark.} The Bahadur property (Definition~\ref{def:bahadur}) is the
only property of the estimator used in the asymptotic analysis. It is worth
stating the main theorems for an abstract estimator satisfying
Definition~\ref{def:bahadur}, then deriving the concrete case from
Lemma~\ref{lem:estimator_bahadur}.

\section{The conditions}
\label{sec:model_conditions}

\begin{definition}[Condition A: second moments]\label{def:condA}
\uses{def:responses}
Condition~\textbf{A} holds if for every arm $k \in [K]$ the response satisfies
$\E|\xi_{1,k}|^2 < \infty$ (equivalently $V_k < \infty$).
\end{definition}

\begin{definition}[Attainable parameter values]\label{def:attainable}
\lean{AlphaRAR.attainableSet}\leanok
\uses{def:estimator}
For each arm $k$, let $V_k^{\mathrm{att}}$ be the set of all values attainable by
the estimator $\hat{\theta}_{n,k}$, over all $n$ and all realizations of the
responses $\{\xi_{j,k}\}$, and set $I_k := \{\theta_k\} \cup V_k^{\mathrm{att}}$.
(The formalized \texttt{AlphaRAR.attainableSet} is the \emph{closure} of the range of
$\hat{\theta}_{\cdot,k}$ over all times and outcomes; being closed it already contains every a.s.\
limit $z_k$ — in particular $\theta_k$ on the infinitely-sampled event — so the explicit union with
$\{\theta_k\}$ is unnecessary.)
\end{definition}

\begin{definition}[Condition B: regular non-sparse target]\label{def:condB}
\uses{def:target, def:attainable}
Condition~\textbf{B} holds if the domain $\hilb$ is open and contains
$I_1 \times \dots \times I_K$, and moreover
\begin{itemize}
  \item $\Target$ is twice differentiable on $\hilb$;
  \item for every $z \in I_1 \times \dots \times I_K$, $\Target(z) \in (0,1)^K$.
\end{itemize}
In particular the target allocation $v = \Target(\Params)$ is non-sparse:
$v_k > 0$ for all $k$.
\end{definition}

\begin{definition}[Sparse target allocation]\label{def:sparse}
\uses{def:target}
A target function $\Target$ is \emph{sparse at $\Params$} if some component of
$v = \Target(\Params)$ is zero.
\end{definition}

\section{Stochastic orders}
\label{sec:model_orders}

The rates in the paper are stated with the probabilistic Landau notation
$o_p$ and $O_p$. We record their definitions and the two facts about them that
are used repeatedly.

\begin{definition}[$o_p$ and $O_p$]\label{def:op_Op}
\lean{AlphaRAR.IsLittleOpOne, AlphaRAR.IsBigOpOne}\leanok
Let $(a_n)_{n\ge1}$ be positive reals and $(X_n)_{n\ge1}$ random variables.
\begin{itemize}
  \item $X_n = o_p(a_n)$ if for all $\varepsilon > 0$,
        $\Prob\!\big(|X_n|/a_n > \varepsilon\big) \to 0$.
  \item $X_n = O_p(a_n)$ if for all $\varepsilon > 0$ there is $M > 0$ with
        $\limsup_n \Prob\!\big(|X_n|/a_n > M\big) \le \varepsilon$
        (boundedness in probability of $X_n / a_n$).
\end{itemize}
\end{definition}

\textbf{Lean remark.} These are exactly the notions
\texttt{Filter.Tendsto} to $0$ in probability and
tightness/boundedness-in-probability. Mathlib already has
\texttt{MeasureTheory} machinery for convergence in probability; a small
API around $o_p/O_p$ (sum, product, scalar multiple, comparison) should be
built first as it is used everywhere.

\begin{lemma}[$o_p$ and $O_p$ under scaling]\label{lem:op_smul}
\uses{def:op_Op}
If $X_n = o_p(a_n)$ and $b_n / a_n \to c \in \R$, then $X_n = o_p(b_n)$
whenever $c \neq 0$; similarly for $O_p$. Moreover $o_p(a_n) + o_p(a_n) = o_p(a_n)$,
$O_p(a_n) + O_p(a_n) = O_p(a_n)$, and $o_p(a_n) + O_p(a_n) = O_p(a_n)$.
\end{lemma}

\begin{proof}
Direct from the definitions and a union bound.
\end{proof}
